% Pursuing Stacks (cotes 134-2 to 134-8): archive cote 134-3, site page 70 (PDF page 71), typescript page 70.
% typescript. Draft compared: ps2.tex:3616-3680.
% AI-transcribed by Claude Opus 5.5 (claude-opus-5-5, adaptive thinking, effort high), Message Batches API, batch ps1 (msgbatch_01Rv5hCaNTtUt4JHx1yqoHPL), 2026-11-12; not yet checked by eye.
\scanpage{134-3/70}
\mspage{\hand{70}}

(We get the previous result by taking the weak equivalences $A\to e_{(Cat)}$
and $id_C : C \to C$.)  To prove the proposition,  we are immediately reduced
to the case when one of the two functors is \xout{xhx} an identity functor\unsure{k}  i.e.
proving the
\uline{Corollary 1}  If  $f$:  $A'\to A$ is a weak equivalence in (Cat),  then for any
\over{\unsure{E}}{C} in (Cat)  $f\hand{x}id_{\hand{C}}$ : $A'xC \to AxC$  is a weak equivalence.

For proving this,  we view the functor $fxid_C$  as a morphism of catego-%
ries ``over C'',  which corresponds to a situation of a morphism of topoi
over a third one
\diagram{\begin{tikzcd}[column sep=tiny]
X' \ar[rr,"f"] \ar[dr] & & X \ar[dl] \\
& Y &
\end{tikzcd}}
\qquad ,

we will say that $f$ is a \uline{weak equivalence} \uline{relative to} $Y$ ,  if not only this
is a weak equivalence,  but remain\hand{s} so by any base change by a localization
functor
\[ Y_{/U} \to U \]
giving rise to
\[ f_{/U} : X'_{/U} \to X_{/U} \qquad . \]
As usual,  standar\over{t}{d} arguments prove that it is enough to take $U$ in a set
of generators of the topos $Y$ .  In case $Y$ =  $C^{\wedge}$\hand{$\hat{}$} ,  we may take $U$ in $C$ .
In case moreover  $X,X'$ are defined by small categories $P,P'$ and a functor
of categories over $C$
\diagram{\begin{tikzcd}[column sep=tiny]
P' \ar[rr,"F"] \ar[dr] & & P \ar[dl] \\
& C &
\end{tikzcd}}
\qquad ,

this condition amounts to demanding that for any $c$ in $C$,  the induced func-%
tor
\[ F_{/c} : P'_{/c} \to P_{/c} \]
be a weak equivalence.  In our case $P=AxC,P'=A'xC$, $F = fxid_C$ ,  the indu-%
ced functor can be identified with  $fx\,id_{C_{/c}}$  .  This reduces us,  for
proving the corollary, to the case when $C$ has a final element.  But consi-%
der now the commutative diagram
\diagram{\begin{tikzcd}
A'xC \ar[r,"\hand{fxid_c}"] \ar[d] & A\deleted{'}xC \ar[d] \\
A' \ar[r,"\hand{f}"] & A
\end{tikzcd}}
\qquad ,
where the vertical arrows are the projections and hence,  by what was
proved before,  weak equivalences.  As $f$ is a weak equivalence,  it follows
that $fxid_C$ is a weak equivalence too,  qed.

T\inserted{h}e proposition above goes somewhat in the direction of looking at
``homotopy properties of (Cat)''  \deleted{\illegible{1}} and ``how far (Cat) is from being a
closed model category in Quillen's sense''.  It is very suggestive \hand{for}
